\section{Limitations}
\label{sec:discussion}

\paragraph{The sizing benchmark.} The stopping rule \eqref{eq:stop} is not anytime-valid. Testing
after every seed and stopping on the first crossing inflates the error of a fixed-sample test, so
its nominal $z$ carries no coverage guarantee and we claim none. What we report is the mistake rate
the rule realises, measured against an ordering fixed on seeds no policy may draw. Every policy is
graded the same way and separated at a common realised rate, so the comparison does not rest on the
boundary being calibrated; a sequentially valid boundary would move all of them together, which we
expect but have not verified. Pairs whose ordering is ambiguous are excluded, so the benchmark says nothing about
near-ties, where the honest answer is to keep training whatever rule is in use.

\paragraph{Scale.} The largest model is a $4$-bit $7$B with $\numAdamScaleSeeds$ seeds, on one
laptop, and the exact tier where the forecast is validated is the furthest from a frontier run.

\paragraph{What the forecast costs.} Estimating the injected term from the eight prompts a run
trains on gives two forecasts $\numModeLargeRepeatRatio\times$ apart, because only about a tenth of
groups are not unanimous. Raising that to $\numAdamScaleWidePrompts$ prompts brings disjoint halves
to within $\numAdamScaleHalfRatio\times$ and costs $141$ minutes against the $260$ that training
all $\numAdamScaleSeeds$ seeds costs. A diagnostic consuming half the experiment it replaces is not
yet a diagnostic at this scale.

\paragraph{Claims we withdrew.} A predicted square-root relation between policy divergence and
outcome spread is not supported by our data. Nor is a short memory horizon: in $\numTiltedEarly$ of
$\numMemorySettings$ settings more of the final variance was injected in the first half of the run.
Token-level reshaping, length normalisation and clipping are outside \eqref{eq:family}.
